%! TEX root = main.tex

\documentclass[italian]{HKNdocument}

\newcommand{\R}{\mathbb{R}}
\newcommand{\inc}{\subseteq}

%Ho creato questi comandi personalizzati perchè nel frattempo sono cambiati, questa
% sintassi permette di usarli in italiano senza cambiare tutte le occorrenze nel file
\newenvironment{definizione}{\begin{definition}}{\end{definition}}
\newenvironment{teorema}{\begin{theorem}}{\end{theorem}}
\newenvironment{osservazione}{\begin{observation}}{\end{observation}}
\newenvironment{esercizio}{\begin{exercise}}{\end{exercise}}
\newenvironment{corollario}{\begin{corollary}}{\end{corollary}}

\courseinfo{
    code = 23ACIMQ,
    name = {Analisi matematica II},
    teacher = {Prof.\ Paolo Tilli},
    program_level = bachelor,
    program = {Matematica per l'Ingegneria},
    program_en = {Mathematics for Engineering},
    start_academic_year = 2024,
    semester = 1
}

\author{Tommaso Pignatelli}

\begin{document}

\include{chapters/introduzione}
\include{chapters/calcolo_differenziale}
\include{chapters/integrali_doppi}
\include{chapters/integrali_tripli}
\include{chapters/integrali_curvilinei}
\include{chapters/campi_conservativi}
\include{chapters/Green}
\include{chapters/integrali_di_superficie}
\include{chapters/Gauss_Green}
\include{chapters/Stokes}
\include{chapters/serie_numeriche}
\include{chapters/serie_di_potenze}
\include{chapters/serie_di_fourier}

\end{document}
